\documentclass{article}
\usepackage{../math_template}

\author{Jonathan Kasongo}
\title{British Math Olympiad 2019 Round 1 --- P5/6}

\begin{document}
\maketitle

\begin{problem}{British Math Olympiad 2019 Round 1 --- P5/6}
Six children are evenly spaced around a circular table. Initially, one has a pile
of $n > 0$ sweets in front of them, and the others have nothing. If a child has
at least four sweets in front of them, they may perform the following move:
eat one sweet and give one sweet to each of their immediate neighbours and
to the child directly opposite them. An arrangement is called perfect if there
is a sequence of moves which results in each child having the same number
of sweets in front of them. For which values of $n$ is the initial arrangement
perfect?
\end{problem}

\begin{solution}{Solution}
Suppose that, without loss of generality, the leftmost child starts off
with the $n$ sweets. Now suppose that the child performs the given
operation a total of $k$ times, the distribution of sweets will look like
so,

\begin{center}
\begin{asy}
size(175);

// Radius of small circles
real r = 0.35;

// Radius of the hexagon
real R = 1.4;

// Draw 6 circles at hexagon vertices
for (int i = 0; i < 6; ++i) {
    real angle = pi/3 * i; // 60° increments
    pair pos = R * dir(degrees(angle)); // position on hexagon
    draw(circle(pos, r));

    // If this is the leftmost circle, add a label on top
    if (i == 3) {
        label("$n - 4k$", pos + (0, r + 0.05), N); // label above the circle
    }

    if (i == 0 || i == 2 || i == 4) {
        label("$k$", pos + (0, r + 0.05), N);
    }

}
\end{asy}
\end{center}

Now suppose each of the children who currently have $k$ sweets now perform
the given operation $a, b$ and $c$ times respectively. We end up with the
following situation,

\begin{center}
\begin{asy}
size(175);

// Radius of small circles
real r = 0.35;

// Radius of the hexagon
real R = 1.5;

// Draw 6 circles at hexagon vertices
for (int i = 0; i < 6; ++i) {
    real angle = pi/3 * i; // 60° increments
    pair pos = R * dir(degrees(angle)); // position on hexagon
    draw(circle(pos, r));

    // If this is the leftmost circle, add a label on top
    if (i == 3) {
        label("$n - 4k + a + b + c$", pos + (0, r + 0.05), N); // label above the circle
    }

    if (i == 4) label("$k-4a$", pos + (0, r+0.05), N);
    if (i == 0) label("$k-4b$", pos + (0, r+0.05), N);
    if (i == 2) label("$k-4c$", pos + (0, r+0.05), N);

    if (i == 5) label("$a+b+c$", pos + (0, r+0.05), N);
    if (i == 1) label("$a+b+c$", pos + (0, r+0.05), N);
}
\end{asy}
\end{center}

Now if the two children who have $a + b + c$ sweets perform the given
operation then the number of sweets the leftmost child has does not change.
Also the number of sweets the children with $k - 4x$ sweets will gain
is the same for all of them. Thus we must have
$k-4a=k-4b=k-4c$ and so $a=b=c$. Also $n - 4k + 3a = k - 4a + y$,
where $y$ is how many sweets the children with $k-4x$ sweets will gain.
\end{solution}

\end{document}
